\documentclass[exercice]{thedocument}%
\startdatakeys
\source{}
\BO{2026}
\cycle{4}
\competencesDuSocle{CA,RA}
\objectifsApprentissage{C4A2db,C4A2bb,C4A2ab}
\automatismes{C4A2ab,C4A3bc}
\enddatakeys
\begin{document}%
Pour chaque affirmation, dire si elle est vraie ou fausse en justifiant~:
\begin{enumerate}
    \item Un nombre premier peut avoir 0 comme chiffre des unités.
    \item Deux nombres premiers peuvent être consécutifs.
    \item Tous les nombres impairs sont premiers.
    \item La somme de deux nombre premier est toujours un nombre premier.
\end{enumerate}
\end{document}
